% Part: lambda-calculus
% Chapter: representability
% Section: currying

\documentclass[../../../include/open-logic-section]{subfiles}

\begin{document}

\olfileid[id]{lam}{rep}{cur}
\olsection{Kurisasi}

Suatu abstrak-$\lambd$, $\lambd[x][M]$, merepresentasikan fungsi satu
argumen. Hal ini cukup membatasi ketika kita ingin mendefinisikan fungsi yang
menerima beberapa argumen. Salah satu cara mengatasinya ialah memperluas
kalkulus-$\lambd$ agar memungkinkan pembentukan pasangan, tripel, dan
seterusnya; dalam hal itu, misalnya, fungsi tiga-tempat $\lambd[x][M]$ akan
mengharapkan argumennya berupa tripel. Namun, lebih praktis jika hal ini
dilakukan melalui \emph{kurisasi} (currying).

Mari kita perhatikan sebuah contoh. Untuk sementara, anggaplah bahwa kita
mempunyai operasi~$+$ dalam kalkulus-$\lambd$. Fungsi penjumlahan beraritas~2,
yakni menerima dua argumen. Namun, abstrak-$\lambd$ hanya memberi kita fungsi
satu argumen: sintaksisnya tidak mengizinkan ekspresi seperti
$\lambd[(x,y)][(x+y)]$. Meskipun demikian, kita dapat mempertimbangkan fungsi
satu-tempat~$f_x(y)$ yang diberikan oleh $\lambd[y][(x+y)]$, yang menambahkan
$x$ pada argumen tunggalnya~$y$. Sebenarnya, ini bukan satu fungsi, melainkan
keluarga fungsi ``tambahkan~$x$'' yang berbeda-beda, satu untuk setiap
bilangan~$x$. Sekarang kita dapat mendefinisikan fungsi satu-tempat lain~$g$
sebagai $\lambd[x][f_x]$. Ketika diterapkan pada argumen~$x$, $g(x)$
mengembalikan fungsi~$f_x$---jadi, nilai-nilainya sendiri adalah fungsi. Jika
kita menerapkan $g$ pada~$x$, kemudian menerapkan hasilnya pada~$y$, kita
memperoleh $(g(x))y=f_x(y)=x+y$. Dengan cara ini, fungsi satu-tempat~$g$ dapat
melaksanakan pekerjaan yang sama dengan fungsi penjumlahan dua-tempat.
``Kurisasi'' sekadar merujuk pada kiat mengubah fungsi dua-tempat menjadi
fungsi satu-tempat yang nilai-nilainya juga fungsi satu-tempat.

Berikut adalah contoh yang sungguh berada dalam sintaksis kalkulus-$\lambd$.
Bagaimana kita merepresentasikan fungsi $f(x,y)=x$? Jika kita ingin
mendefinisikan fungsi yang menerima dua argumen dan mengembalikan argumen
pertama, kita dapat menulis $\lambd[x][\lambd[y][x]]$, yang secara harfiah
merupakan fungsi yang menerima argumen~$x$ dan mengembalikan fungsi
$\lambd[y][x]$. Fungsi $\lambd[y][x]$ menerima argumen lain~$y$, tetapi
mengabaikannya dan selalu mengembalikan~$x$. Mari kita lihat apa yang terjadi
ketika $\lambd[x][\lambd[y][x]]$ diterapkan pada dua argumen:
\begin{align*}
  (\lambd[x][\lambd[y][x]])MN
  \bredone & (\lambd[y][M])N \\
  \bredone & M
\end{align*}

Secara umum, untuk menulis fungsi dengan parameter $x_1$, \dots,~$x_n$ yang
didefinisikan oleh suatu suku~$N$, kita dapat menulis
$\lambd[x_1][\lambd[x_2][\ldots\lambd[x_n][N]]]$. Jika kita menerapkan $n$
argumen padanya, kita memperoleh:
\begin{multline*}
  (\lambd[x_1][\lambd[x_2][\ldots\lambd[x_n][N]]]) M_1 \dots M_n \bredone\\
  \begin{aligned}
  \bredone {} & (\Subst{(\lambd[x_2][\ldots\lambd[x_n][N]])}{M_1}{x_1}) M_2
  \dots M_n\\
   \eqs {} & (\lambd[x_2][\ldots\lambd[x_n][\Subst{N}{M_1}{x_1}]]) M_2
            \dots M_n \\
  \vdots & \\
  \bredone {} &\Subst{\Subst{N}{M_1}{x_1}\ldots}{M_n}{x_n}
  \end{aligned}
\end{multline*}
Baris terakhir secara harfiah berarti menyubstitusikan $M_i$ sebagai pengganti
$x_i$ dalam badan definisi fungsi, tepat seperti yang kita inginkan ketika
menerapkan beberapa argumen pada suatu fungsi.
\end{document}
